\section{Trabajo Relacionado}

\subsection{Interpolación Clásica de EEG}
El problema de interpolar canales EEG faltantes ha sido históricamente abordado utilizando métodos geométricos y estadísticos. La línea base más ampliamente adoptada es la interpolación por Spherical Spline \cite{perrin1989}, que modela el cuero cabelludo como una esfera continua y proyecta potenciales observados a ubicaciones de electrodos faltantes utilizando polinomios ponderados por distancia. Otras variantes incluyen Interpolación de Función de Base Radial (RBFI) \cite{jager2016} y Thin-Plate Splines. Aunque estos métodos son computacionalmente económicos y simples de implementar, operan estrictamente en el dominio espacial para muestras de tiempo aisladas. Asumen inherentemente que los potenciales eléctricos varían suavemente en la superficie del cuero cabelludo, pero ignoran completamente la conectividad funcional subyacente y la continuidad temporal de los generadores neurofisiológicos.

Como alternativa, métodos estadísticos como Análisis de Componentes Independientes (ICA) intentan proyectar los datos a un conjunto de fuentes independientes, reconstruir los sensores faltantes a partir de las señales de fuentes, y proyectar de vuelta al espacio de sensores. Sin embargo, ICA tiene dificultades con artefactos no estacionarios y frecuentemente suprime ráfagas localizadas de alta frecuencia al reconstruir canales faltantes.

\subsection{Procesamiento de Señales en Grafos (GSP)}
La interpolación basada en grafos para datos faltantes en señales estructuradas representa un cambio significativo de métricas de distancia puramente euclidianas. Las formulaciones tempranas de Narang y Ortega \cite{narang2013} demostraron que la interpolación puede formularse como un problema de optimización donde las señales se restringen a ser suaves respecto a una topología de grafo. Trabajo teórico adicional de Puy \emph{et al.} \cite{puy2018} aclaró las condiciones bajo las cuales las señales de grafo de banda limitada pueden ser recuperadas con precisión a partir de observaciones parciales.

El marco GSP general, como es detallado exhaustivamente por Shuman \emph{et al.} \cite{shuman2013} y Ortega \emph{et al.} \cite{ortega2018}, establece la intuición de que un Laplaciano de grafo puede reemplazar al operador de derivada espacial tradicional cuando los datos residen en una variedad irregular, como el cuero cabelludo humano. Métodos más recientes han expandido esto integrando Espacios de Hilbert de Núcleo Reproductor (RKHS) para recuperación continua, como el algoritmo BGSRP \cite{zhang2024bgsrp}.

Los avances teóricos recientes fortalecen aún más esta línea de trabajo proporcionando garantías explícitas de recuperación y límites de coherencia para recuperación de señales de grafo dispersas \cite{morgenstern_theoretical_2025}, mientras que estudios aplicados en procesamiento de señales biomédicas han mostrado que las representaciones espectrales de grafos pueden mejorar la robustez bajo condiciones neurofisiológicas no estacionarias \cite{miri_spectral_2024}. En conjunto, estos desarrollos motivan evaluar métodos de interpolación no solo por error de amplitud, sino también por fidelidad temporal y espectral.

\subsection{Regularización Temporal y la Brecha de Benchmarking}
Mientras que los métodos GSP estáticos incorporan exitosamente la topología espacial, aún reconstruyen señales fotograma por fotograma. Los avances recientes en señales de grafo variables en el tiempo \cite{giraldo2022} abordan esto introduciendo restricciones de suavidad de Sobolev en el eje temporal, optimizando conjuntamente derivadas espaciales y temporales.

Esta dirección se alinea con revisiones más amplias sobre procesamiento de señales de grafo variables en el tiempo, que enfatizan que el modelado temporal explícito es crucial siempre que las observaciones estructuradas en grafos evolucionan dinámicamente \cite{yan_signal_2026, krishnan_learning_2026}.

Simultáneamente, evidencia práctica de aplicaciones de EEG e imputación de sensores indica que la interpolación consciente de grafos mejora consistentemente la robustez sobre métodos puramente geométricos, especialmente cuando la pérdida de canales es estructurada o severa \cite{weinstein_imputing_2023, jiang_graph-based_2021}. Estudios complementarios de interpolación de grafos no suave también advierten que el sobre-suavizado puede ocultar transitorios locales, reforzando la necesidad de regularización espaciotemporal equilibrada \cite{mazarguil_non-smooth_2022}.

A pesar de la proliferación de estos algoritmos avanzados, la literatura sufre de una brecha de benchmarking sustancial. Los métodos novedosos son frecuentemente evaluados contra líneas base clásicas (p. ej., Spherical Splines) ejecutadas utilizando sus configuraciones de software por defecto, en lugar de ser sintonizadas rigurosamente para el conjunto de datos específico. Además, las evaluaciones a menudo se restringen a un único conjunto de datos o protocolo de enmascaramiento, obscureciendo si un algoritmo es universalmente robusto o simplemente sobreadaptado a un setup experimental específico.

El presente trabajo aborda esta brecha directamente. En lugar de proponer una teoría de interpolación nueva, introducimos un pipeline de benchmarking comprehensivo y multi-dataset impulsado por optimización sistemática de hiperparámetros (vía Optuna). Esto asegura que las líneas base clásicas y los métodos modernos de GSP temporal sean comparados en sus techos teóricos absolutos bajo condiciones de enmascaramiento diversas y realistas.
